\section{Warm starting reduces repeated nonlinear work}
\label{sec:background}

\subsection{Transient simulation exposes a stage-level target}

After time discretization, a circuit simulator solves a nonlinear modified nodal equation at time $t_n$,
\begin{equation}
F_n(x_n;x_{n-1},u_n)=0,
\label{eq:nonlinear}
\end{equation}
where $x_n$ contains node voltages and branch variables, $x_{n-1}$ carries integration history, and $u_n$ contains independent sources. Newton iteration computes
\begin{equation}
J_n(x_n^{(k)})\Delta x_n^{(k)}=-F_n(x_n^{(k)}), \quad
x_n^{(k+1)}=x_n^{(k)}+\alpha_k\Delta x_n^{(k)}.
\label{eq:newton}
\end{equation}
The Jacobian $J_n$ and damping coefficient $\alpha_k$ follow the native solver. Continuation methods can create several nonlinear stages at one time point by changing $g_{\min}$ or source strength~\cite{kundert1990steady}.

The native solver supplies an initial state $x^{\mathrm{base}}$ at each stage. \system{} predicts a correction $\widehat{\delta x}$ and forms
\begin{equation}
x^{\mathrm{warm}}=x^{\mathrm{base}}+\widehat{\delta x}.
\label{eq:warm}
\end{equation}
The training target uses the converged state from that same stage,
\begin{equation}
\delta x^*=x^{\mathrm{conv}}-x^{\mathrm{base}}.
\label{eq:target}
\end{equation}
This target aligns supervision with the deployment goal. One-step Newton prediction would imitate an intermediate trajectory that the warm start seeks to shorten.

\subsection{Overlapping cells create proposal conflicts}

Let the transistor netlist form a graph $G=(V,E)$. Each standard-cell instance $c_i$ induces a local subgraph $G_i$ with internal variables $x_i^{\mathrm{int}}$ and port variables $x_i^{\mathrm{port}}$. Internal variables belong to one instance. A shared net can appear as an output pin, several input pins, and an interconnect node across multiple instances. The local subgraphs therefore overlap on ports.

For a cell type $\tau$, all instances share proposer parameters $\theta_\tau$. A local-to-global map $\pi(i,j)=v$ assigns local variable $j$ of instance $i$ to global variable $v$. Proposers create the candidate set
\begin{equation}
\mathcal{P}_v=\{(p_{i,j},\tau_i,r_{i,j})\mid \pi(i,j)=v\},
\label{eq:candidates}
\end{equation}
where $r_{i,j}$ records the pin role and instance metadata. Candidate disagreement carries physical information. A driving output and several input loads have different authority over a shared net, and their local views can disagree most during switching.

\subsection{Strict success couples speed and accuracy}

We count a paired stage as a strict success when the baseline and \system{} complete the same simulation task, satisfy the same native convergence tests, and meet the waveform tolerance. Every fallback remains visible in the record. For baseline and warm-start Newton counts $K_b$ and $K_w$, we report iteration reduction $(K_b-K_w)/K_b$. For end-to-end times $T_b$ and $T_w$, we report speedup $T_b/T_w$. $T_w$ includes input construction, model inference, guards, restarts, and native solving.

Residual equality alone cannot establish waveform equality when nonlinear equations have several attraction regions. We therefore compare the complete output waveform, final nonlinear residual, and Kirchhoff current-law residual. This contract turns learned prediction error into a performance concern while the native solver retains responsibility for accepted accuracy.

